\section{Module 1: Classical Machine Learning}

\begin{frame}{Why Classical Machine Learning Comes First}
\begin{itemize}
  \item Classical methods (k-NN, PCA, clustering, decision trees) have clear geometric or statistical intuition and can be implemented from scratch with NumPy.
  \pause
  \item The same core ideas -- distance, variance, loss, generalization -- reappear inside every model taught later in the course.
  \pause
  \item This module builds the practical and geometric vocabulary that every later module assumes.
\end{itemize}
\pause
\begin{block}{Course principle}
We start from methods that are simple enough to analyze carefully, then revisit the same questions for stronger and more complex models.
\end{block}
\end{frame}

\begin{frame}{Reading Resources}
\begin{itemize}
  \item We do not follow a single textbook strictly.
  \pause
  \item Core theory reference: \textit{Foundations of Machine Learning} (Mohri et al.) -- learning theory and generalization guarantees.
  \pause
  \item Statistical learning reference: \textit{The Elements of Statistical Learning} (Hastie, Tibshirani, Friedman) -- trees, boosting, ensembles.
  \pause
  \item Probabilistic modeling reference: \textit{Pattern Recognition and Machine Learning} (Bishop) -- mixture models, EM.
  \pause
  \item Practical reference: scikit-learn docs (\url{https://scikit-learn.org/stable/}), used throughout to sanity-check from-scratch NumPy implementations.
  \pause
  \item We also use selected papers, lecture notes, benchmark datasets, and practical case studies.
\end{itemize}
\end{frame}
